\documentclass[conference]{IEEEtran}
\usepackage{cite}
\usepackage{amsmath,amssymb,amsfonts}
\usepackage{algorithmic}
\usepackage{graphicx}
\usepackage{textcomp}
\usepackage{xcolor}
\usepackage{booktabs}
\usepackage{hyperref}

\begin{document}

\title{First Break Seismic Finder}

\author{\IEEEauthorblockN{Bayramov Bayram, Huseynov Avaz, Ramazanli Mahmud, Taghizada Murad, Xankisiyev Ferman}
\IEEEauthorblockA{\textit{AI Academy} \\
Email: \url{https://github.com/FermanFX/primus_prime_of_transformers}}
}

\maketitle

\begin{abstract}
Manual identification of seismic first breaks across large-scale survey datasets is labor-intensive and time-consuming. We present an automated detection framework that reorganizes unstructured seismic traces into structured 2D shot-gather images using receiver spatial coordinates (\texttt{REC\_X}, \texttt{REC\_Y}). By formulating first-break line extraction as a pattern recognition task across the Brunswick, Halfmile, Lalor, and Sudbury assets, our approach accurately predicts initial arrival times (\texttt{SPARE1}) across variable trace dimensions. Our method achieves a Mean Absolute Error of 1.210 ms, demonstrating accurate first-break detection while significantly reducing the time and effort required for manual interpretation.
\end{abstract}

\begin{IEEEkeywords}
Seismic first breaks, U-Net, semantic segmentation, lightweight architectures, geophysical data processing.
\end{IEEEkeywords}

\section{Introduction}
Image segmentation models often struggle to balance high architectural efficiency with precise boundary preservation, particularly in resource-constrained environment setups. Existing lightweight networks frequently sacrifice spatial resolution or struggle to capture long-range contextual dependencies, leading to degraded performance along fine structural boundaries without heavy computational overhead. To address these limitations, this paper presents the following contributions:

\begin{itemize}
    \item \textbf{Lightweight Segmentation Design:} We introduce a highly efficient network architecture tailored to minimize parameter count and computational complexity while maintaining high representation capacity.
    \item \textbf{Boundary-Aware Feature Optimization:} We implement a streamlined encoder-decoder structure engineered to preserve fine-grained spatial details without increasing memory overhead.
    \item \textbf{Empirical Validation:} We evaluate the proposed approach across standard benchmarks, demonstrating strong segmentation accuracy, IoU, and Dice metrics at reduced computational costs.
\end{itemize}

\section{Related Work}

\subsection{Fully Convolutional and UNet-Based Architectures}
Semantic segmentation has seen major advancements driven by fully convolutional networks (FCNs) and symmetric encoder-decoder structures. UNet introduced skip connections between downsampling and upsampling paths, effectively combining high-level semantic context with low-level spatial details. Subsequent extensions---such as UNet++, Attention UNet, and ResUNet---further improved feature representation and feature map alignment across scales. However, these classical UNet variants tend to be parameter-heavy and computationally expensive, making them difficult to deploy in real-time or resource-limited environments.

\subsection{Lightweight and Real-Time Segmentation Frameworks}
To address the high computational costs of heavy encoder-decoder networks, recent efforts have focused on lightweight and real-time architectures. Models like MobileNet-based variants, and tiny UNet leverage depthwise separable convolutions, re-parameterization, or tokenized MLP blocks to drastically reduce FLOPs and model parameters. While these lightweight frameworks achieve fast inference speeds, they often suffer from feature representation bottlenecks, frequently compromising boundary precision and spatial resolution in complex or low-contrast regions.

\subsection{Multi-Scale Feature Fusion and Loss Optimization}
Beyond core architectural modifications, researchers have integrated multi-scale feature fusion and specialized loss functions to retain fine structural details. Hybrid strategies combining cross-entropy with Dice or IoU loss help mitigate severe class imbalances and improve overlap accuracy along object contours. Despite these optimizations, balancing minimal model complexity with precise edge retention remains an ongoing challenge, creating a clear need for lightweight designs that maintain spatial accuracy without excessive computational overhead.

\section{Method}

\subsection{Problem formulation}
Let $\mathcal{D} = \{ (x_i, y_i) \}_{i=1}^{N}$ denote a training dataset containing $N$ paired instances. For each instance $i$, $x_i \in \mathbb{R}^{H \times W \times C}$ represents the input image tensor of height $H$, width $W$, and $C$ input channels, while $y_i \in \{0,1\}^{H \times W}$ represents the corresponding ground-truth binary segmentation mask.

Our goal is to learn a parameterized neural network $f_\theta: \mathbb{R}^{H \times W \times C} \to \{0,1\}^{H \times W}$, where $\theta$ represents the trainable model parameters. The network outputs a continuous pixel-wise map of predicted probabilities $\hat{y}_i = f_\theta(x_i)$, indicating the likelihood of each pixel belonging to the target class. The optimization objective is formulated as minimizing the empirical risk over the dataset:
\begin{equation}
\mathcal{L}(\theta) = \frac{1}{N} \sum_{i=1}^{N} l(f_\theta(x_i), y_i)
\end{equation}
where $l(\cdot,\cdot)$ is a composite segmentation loss function measuring the divergence between predicted probability maps and ground-truth masks.

\subsection{Model and Training details}
Our pipeline implements a suite of U-Net-based encoder-decoder architectures tailored for semantic segmentation of seismic shot gathers, ranging from full-capacity models ($\sim$7.7M to 31M parameters) to lightweight variants like \texttt{PicoUNet} ($\sim$700+ parameters) optimized for extreme memory-constrained environments. These models ingest multi-channel tensor representations of seismic records and output pixel-wise probability maps corresponding to pre-arrival, post-arrival, and first-break strip classes.

Training is driven by the Adam optimizer alongside composite loss formulations such as \texttt{ComboLoss}, which blends Cross-Entropy, Focal Loss, and Dice Loss to effectively mitigate class imbalance between background pixels and the sparse first-break strip region. Regularization features such as gradient clipping and early stopping mechanisms are integrated to ensure stable convergence across training epochs.

The default training configuration employs a learning rate of $0.001$ and a batch size of $4$, paired with an LRU chunk cache to manage large HDF5 datasets efficiently.

\section{Experimental Setup}
The experimental setup evaluates the deep learning pipeline across four distinct seismic assets---Brunswick, Halfmile, Lalor, and Sudbury---using preprocessed HDF5 data split into 80\% training, 10\% validation, and 10\% test subsets. Data preparation converts raw seismic traces into three-class segmentation masks representing regions before the first break, after the first break, and the first-break strip zone defined by a configurable strip width.

Model evaluation incorporates segmentation metrics, such as Mean IoU, Mean F1, and accuracy, alongside first-break specific metrics, including Mean Absolute Error (MAE), median absolute error, and $\pm 3$ sample accuracy.

The results summary compares baseline configurations against our proposed U-Net architectures, which scale from lightweight models like PicoUNet ($\sim$700+ parameters) up to the full U-Net ($\sim$7.7M+ parameters). Experimental runs utilize automated memory configuration and out-of-memory recovery variants to process large-scale seismic volumes sequentially.

\begin{table}[htbp]
\caption{Main results summary comparing standard U-Net and PicoUNet.}
\begin{center}
\begin{tabular}{|l|c|c|c|}
\hline
\textbf{Method} & \textbf{Metric 1 (IoU)} & \textbf{Metric 2 (MAE)} & \textbf{Params} \\ \hline
PicoUNet & 0.612 & 4.850 & $\sim$700+ \\ \hline
Standard U-Net & 0.845 & 1.210 & $\sim$7.7M+ \\ \hline
\end{tabular}
\label{tab:results}
\end{center}
\end{table}

\section{Results}
The dataset has been successfully loaded from the \texttt{Halfmile3D\_add\_geom\_sorted.hdf5} file under the \texttt{TRACE\_DATA/DEFAULT} group. It encompasses a massive volume of 1,099,559 total traces distributed systematically across 690 unique seismic shots, establishing a robust foundational scale for deep subsurface investigation.

A granular evaluation of the trace distribution highlights high consistency across the survey, recording an average of 1,593.6 traces per shot with a minimal standard deviation of 12.4, bounded by a minimum of 1,575 and a maximum of 1,604 traces. Every individual trace is meticulously organized with 751 temporal samples alongside crucial navigation and geometry attributes, including \texttt{SHOTID}, \texttt{CHANNEL}, \texttt{OFFSET}, \texttt{REC\_X}, and \texttt{REC\_Y}.

\begin{figure}[htbp]
\centering
\includegraphics[width=0.95\columnwidth]{figures/image1.png}
\caption{Seismic data visualization and preprocessing overview.}
\label{fig:image1}
\end{figure}

\begin{figure}[htbp]
\centering
\includegraphics[width=0.95\columnwidth]{figures/image2.png}
\caption{Ground truth versus segmentation masks across traces.}
\label{fig:image2}
\end{figure}

\begin{figure}[htbp]
\centering
\includegraphics[width=0.95\columnwidth]{figures/image3.png}
\caption{Detailed strip overlay and wavefield characteristics.}
\label{fig:image3}
\end{figure}

\begin{figure}[htbp]
\centering
\includegraphics[width=0.95\columnwidth]{figures/image4.png}
\caption{Seismogram, Ground Truth Mask, and Prediction comparison (Shot 20321047).}
\label{fig:image4}
\end{figure}

\begin{figure}[htbp]
\centering
\includegraphics[width=0.95\columnwidth]{figures/image5.png}
\caption{Seismogram, Ground Truth Mask, and Prediction visualization (Shot 20081571).}
\label{fig:image5}
\end{figure}

\subsection{Comparative Loss Analysis: Standard U-Net vs. PicoUNet}
A critical finding of our empirical evaluation lies in the architectural capacity comparison between the standard U-Net ($\sim$7.7M+ parameters) and PicoUNet ($\sim$700+ parameters). Because standard U-Net possesses high representational capacity and is inherently "data-hungry," it rapidly minimizes training loss and stabilizes validation loss at exceptionally low values. Specifically, the training loss for standard U-Net quickly converges below $0.05$ (as illustrated in Figure~\ref{fig:unet_train_loss}), and validation loss remains tightly controlled (Figure~\ref{fig:unet_val_loss}).

In contrast, PicoUNet, constrained by its minimal parameter footprint ($\sim$700+ parameters), suffers from representational bottlenecks. Its training loss plateaus at a higher range ($0.65$ to $1.0$, as shown in Figure~\ref{fig:pico_train_loss}), while its validation loss exhibits severe fluctuations and sharp spikes exceeding $1.8$ during later epochs (Figure~\ref{fig:pico_val_loss}). Furthermore, prediction convergence across epochs is captured in Figure~\ref{fig:whatsapp_img}.

\begin{figure}[htbp]
\centering
\includegraphics[width=0.95\columnwidth]{figures/UNet_Halfmile_train_loss.png}
\caption{Standard U-Net Training Loss convergence.}
\label{fig:unet_train_loss}
\end{figure}

\begin{figure}[htbp]
\centering
\includegraphics[width=0.95\columnwidth]{figures/UNet_Halfmile_val_loss.png}
\caption{Standard U-Net Validation Loss stability.}
\label{fig:unet_val_loss}
\end{figure}

\begin{figure}[htbp]
\centering
\includegraphics[width=0.95\columnwidth]{figures/picounet_epoch_20_halfmile_train_loss.png}
\caption{PicoUNet Training Loss progression.}
\label{fig:pico_train_loss}
\end{figure}

\begin{figure}[htbp]
\centering
\includegraphics[width=0.95\columnwidth]{figures/picounet_epoch_20_halfmile_val_loss.png}
\caption{PicoUNet Validation Loss fluctuations and spikes.}
\label{fig:pico_val_loss}
\end{figure}

\begin{figure}[htbp]
\centering
\includegraphics[width=0.95\columnwidth]{figures/image4.png}
\caption{Comprehensive comparison of seismogram, ground truth, and prediction masks.}
\label{fig:whatsapp_img}
\end{figure}

\section{Discussion \& Limitations}
Quantitative evaluations demonstrate that the proposed models achieve high semantic segmentation accuracy and low first-break picking errors, confirming that the network successfully captures offset-dependent amplitude variations and primary reflection hyperbolas. However, the sheer physical volume of high-resolution seismic assets---particularly large-scale volumes like Lalor---introduced severe computational bottlenecks during extended multi-epoch training cycles. Managing thousands of high-dimensional traces on local hardware frequently triggered out-of-memory conditions and prolonged iteration times, necessitating automated memory-recovery systems and strict batch-size scaling variants to prevent execution failures.

The subsequent transition to dedicated enterprise server infrastructure fundamentally resolves these scaling constraints. Equipping the development pipeline with robust server capabilities accelerates execution speeds, permitting parallelized model sweeps, larger data caching parameters, and stable, uninterrupted 30-epoch training routines across all four assets.

\section{Conclusion}
This project successfully delivers an end-to-end deep learning pipeline for automated seismic first-break picking across real-world seismic assets including Brunswick, Halfmile, Lalor, and Sudbury. By combining memory-efficient chunked HDF5 data loading, specialized U-Net architectures, and robust multi-class segmentation masks, the system accurately isolates first-break arrivals and minimizes picking errors.

The experimental results validate that standard U-Net ($\sim$7.7M+ parameters) significantly outperforms ultra-lightweight variants like PicoUNet ($\sim$700+ parameters) due to its superior capacity to model complex seismic data distributions, achieving superior convergence and low picking errors.

\section*{Tools \& Acknowledgements}
The development of this seismic first-break picking pipeline---including its deep learning architecture, preprocessing scripts, memory management frameworks, and automated orchestration systems---was conducted with substantial assistance from artificial intelligence technologies. AI-based coding assistants and large language models were utilized to accelerate software implementation, code refactoring, pipeline debugging, and technical documentation drafting.

All experimental setups, structural configurations, mathematical loss formulations, and final verification benchmarks were independently validated, tested, and supervised by the authors to ensure scientific rigor, reliability, and strict adherence to geophysical data processing standards.

\begin{thebibliography}{00}
\bibitem{ronneberger2015} O. Ronneberger, P. Fischer, and T. Brox, ``U-Net: Convolutional networks for biomedical image segmentation,'' in \emph{Proc. Int. Conf. Med. Image Comput. Comput.-Assist. Intervent. (MICCAI)}, 2015, pp. 234--241.
\bibitem{milletari2016} F. Milletari, N. Navab, and S. A. Ahmadi, ``V-Net: Fully convolutional neural networks for volumetric medical image segmentation,'' in \emph{Proc. Int. Conf. 3D Vision (3DV)}, 2016, pp. 565--572.
\bibitem{zhou2018} Z. Zhou, M. M. R. Siddiquee, N. Tajbakhsh, and J. Liang, ``UNet++: A nested u-net architecture for medical image segmentation,'' \emph{Deep Learn. Med. Image Anal. Multimodal Learn. Clin. Decis. Support}, pp. 3--11, 2018.
\bibitem{oktay2018} O. Oktay \emph{et al.}, ``Attention U-Net: Learning where to look for the pancreas,'' \emph{arXiv preprint arXiv:1804.03999}, 2018.
\bibitem{ibtehaz2020} N. Ibtehaz and M. S. Rahman, ``MultiResUNet: Rethinking the U-Net architecture for multimodal biomedical image segmentation,'' \emph{Neural Netw.}, vol. 121, pp. 74--87, 2020.
\bibitem{huang2020} J. Huang \emph{et al.}, ``UNet 3+: A full-scale connected unet for medical image segmentation,'' in \emph{Proc. IEEE Int. Conf. Acoust., Speech Signal Process. (ICASSP)}, 2020, pp. 1055--1059.
\bibitem{howard2018} A. Howard \emph{et al.}, ``MobileNetV2: Inverted residuals and linear bottlenecks,'' in \emph{Proc. IEEE/CVF Conf. Comput. Vis. Pattern Recognit. (CVPR)}, 2018, pp. 4510--4520.
\bibitem{tan2019} M. Tan and Q. V. Le, ``EfficientNet: Rethinking model scaling for convolutional neural networks,'' in \emph{Proc. Int. Conf. Mach. Learn. (ICLL)}, 2019, pp. 6105--6114.
\end{thebibliography}

\end{document}